\section{Sequential Bayesian Meta Analysis for Tracking Scientific Consensus}
\label{Sec:QuantifyingLearning}

We develop a sequential meta-analysis method to track how evidence accumulates as new studies arrive. SMART sequentially updates a posterior distribution of the target parameter after each new study, and we quantify how much the posterior distribution changes at each step. Our method involves sequential updating with a Bayesian meta-analysis model, where each new study revises our collective beliefs, and a measurement of posterior change, which summarizes how those beliefs shift over time.

\subsection{Sequential Bayesian updating}
Consider a target parameter $\vartheta$, such as a causal effect that we are trying to estimate, and suppose that studies $y_t$ will appear sequentially. Sequential Bayesian updating revises our beliefs about $\vartheta$ as new studies appear.  Before observing any studies, we express initial uncertainty about $\vartheta$ with a prior distribution $p(\vartheta)$. After observing study $y_t$, we update the posterior distribution using Bayes' rule,
\begin{align}
  \label{eq:seq_update}
  p(\vartheta \mid y_{1:t})
  \propto p(y_t \mid \vartheta)\,p(\vartheta \mid y_{1:t-1}).
\end{align}
Here $p(y_t \mid \vartheta)$ is the likelihood model of a study, and $p(\vartheta \mid y_{1:t-1})$ is the posterior distribution after studies $1$ through $t-1$. Notice that the posterior after study $t$ becomes the prior before study $t+1$. These sequential updates formally accumulate scientific evidence across multiple studies.

Sequential updating requires a Bayesian meta-analysis model---a joint probability model $p(\vartheta, y_{1:t})$ that connects the parameter $\vartheta$ to the observed data. The posterior distribution, defined in Eq.~\ref{eq:seq_update}, is derived from this joint model. We next describe two hierarchical models for sequential Bayesian meta-analysis: the standard Bayesian random effects model (Section~\ref{sec:random-effects-model}) and an extension, which we call the labeled random effects model (Section~\ref{sec:labeled-random-effects-model}).

\subsection{Bayesian Random Effects for Meta Analysis}
\label{sec:random-effects-model}

The random effects model is a workhorse in the meta-analysis literature~\citep{dersimonian1986meta,Higgins_Thompson_Spiegelhalter_2009,Spence_Steinsaltz_Fanshawe_2016}. It describes studies whose measurements differ systematically from the target parameter $\vartheta$. Formally, study $t$ produces measurement $y_t$ according to the hierarchical model
\[
  y_t \sim \mathcal{N}(\vartheta + b_t, \sigma_t^2), \quad b_t \sim \mathcal{N}(0, \tau^2),
\]
where $b_t$ is the systematic bias of study $t$, $\sigma_t^2$ is the known sampling variance, and $\tau^2$ is the variance of biases across studies. Given a prior on the target parameter, $p(\vartheta)$, we sequentially update our posterior distribution as in Eq.~\ref{eq:seq_update}.

In the random effects model, the bias term $b_t$ can represent two related ideas. \textit{Statistical bias} is a systematic discrepancy between the expected value of an estimate and the true parameter. For example, observational studies can systematically differ from randomized experiments. \textit{Conceptual bias} is a systematic gap between what a study measures and the scientific target of interest. For example, studies on rats measure treatment effects that differ systematically from effects in humans; or studies from limited populations differ systematically from the full population of interest. Mathematically, conceptual and statistical biases are identical---each $b_t$ captures a systematic difference from $\vartheta$.

The random-effects model is appropriate for sequential meta-analysis because it allows our uncertainty about the target parameter to increase or decrease with each new study. A new study that aligns closely with previous findings typically reduces posterior uncertainty. But a surprising or conflicting study can increase posterior uncertainty, as it suggests larger systematic biases than previously believed. This flexibility distinguishes the random-effects model from the simpler fixed-effect model, which assumes no systematic biases and thus cannot increase uncertainty in response to conflicting evidence \citep{Cooper_2009}; see Appendix~\ref{sec:Meta_analysis}.


Though flexible, the random-effects model also makes strong assumptions. Specifically, through its prior distribution, it assumes that the biases
$b_t$ have mean zero. This assumption encodes the belief that studies are unbiased on average. But if some studies systematically differ, for example, due to known methodological biases, then this assumption is violated. (In classical meta-analysis, excluding studies that are known to be flawed is common practice.)

The zero-mean assumption might seem unpalatable, but note that the average bias of the studies is not identifiable from the data alone. Without additional assumptions, the parameter $\vartheta$ cannot be distinguished from the average bias across studies \citep{gustafson2015bayesian}. It is this identification problem that motivates the model of the next section.

\subsection{Labeled Random Effects for Meta Analysis}
\label{sec:labeled-random-effects-model}

Suppose we know specific study characteristics that correlate with biases, for example, study design, population sampled, or methodology. Then we can partially relax the zero-mean assumption by explicitly modeling these characteristics. The \textit{labeled random-effects model} allows systematic biases to depend on observable labels associated with each study. It provides additional structure that can help protect our estimates of the target parameter from systematic biases.

Formally, suppose each study $t$ is associated with a known methodological label $\ell_t$, indicating the methodology or category of the study, and suppose there are $L$ possible methodologies. Each methodology $k$ introduces a bias $\gamma_k$, which we model hierarchically as
\begin{align}
  \gamma_k &\sim \mathcal{N}(0, \kappa_k^2) \quad ( k = 1, \dots, L).
\end{align}
Here, $\kappa_k^2$ describes how much uncertainty we have about the bias associated with methodology $k$. A reliable methodology has a small $\kappa_k^2$, indicating little expected bias, whereas an uncertain or potentially biased methodology has a large $\kappa_k^2$.

With these labeled biases, our model of the observed measurement $y_t$ becomes
\[
  y_t \sim \mathcal{N}(\vartheta + \gamma_{\ell_t} + b_t, \sigma_t^2), \quad b_t \sim \mathcal{N}(0, \tau^2),
\]
where $b_t$ is the study-specific random effect as in the previous model, and $\gamma_{\ell_t}$ is the systematic bias due to the methodology of study $t$. Again, we can update the posterior sequentially with
Eq.~\ref{eq:seq_update}.
Beliefs about biases $\kappa^2_k$ may vary over time. When this is the case, we can quantify how much has been learned retrospectively (i.e., gauging the contribution of a study given our current understanding of biases). Alternatively, we can calculate the posterior at every time point with the $\kappa^2_k$'s describing the beliefs at the time. This gives us a picture of how influential studies were at the time they entered the literature.

This extension explicitly captures method-specific bias, and it allows us to incorporate expert knowledge about which methods are more or less trustworthy. Specifically, the degree to which we believe each methodology to be biased is encoded in $\kappa_{k}$. This, in turn, determines how much the studies that use that methodology contribute to the posterior of $\vartheta$.  The labeled random-effects model requires
that each study is labeled with its methodology, but it improves both
identification and interpretability of the parameter $\vartheta$.

In meta-analysis, studies that are deemed biased are often removed
\citep{Campbell_guidlines}. The labeled random-effects model provides a formal rationale for this often informal practice, down-weighting rather than excluding altogether data from methodologically questionable studies. Note also that it is
different from meta-regression
\citep{Thompson_Higgins_2002,Stanley_2005,stanley2012meta}, which aims
to explain the heterogeneity of reported effects, using study-level covariates.

Unlike the random-effects model, which assumes all studies to be exchangeable, the labeled random-effects model only assumes exchangeability within each label. Benefits of the relaxation of this assumption have been noted in \citep{Raudenbush_1985}.

Finally, we note that both the random-effects and labeled random-effects models rely on substantive prior assumptions. The random-effects model assumes the biases have zero mean; the labeled random-effects model further assumes that each methodology's bias has zero mean and that its variance ($\kappa_k^2$) is substantively informed. One might wonder whether these assumptions can be relaxed further, allowing data to determine the average bias or the variance parameters directly. Unfortunately, such an ``agnostic'' approach faces fundamental identification problems~\citep{gustafson2015bayesian,Gerber_Green_Kaplan_2004}: without prior constraints, the data cannot distinguish the target parameter $\vartheta$ from systematic shifts in bias or uncertainty about biases. Consequently, fully agnostic Bayesian meta-analysis is impossible. Priors---explicitly or implicitly---are unavoidable. The principled solution is sensitivity analysis, which transparently explores how conclusions about $\vartheta$ depend on prior assumptions~\citep{Li_Ding_Mealli_2022,Mikhaeil_2025}.


\subsection{Sequential Measurements of Posterior Change}
\label{Sec:LearningMetric}

Our goal is to sequentially update a collective posterior and measure how much our collective belief about the target parameter $\vartheta$ changes after each new study. Tracking posterior change reveals how evidence accumulation shapes our understanding of the target parameter over time. It can help identify influential studies, and quantify the point at which a scientific consensus begins to stabilize.

Formally, let $p_t \triangleq p(\vartheta \mid y_{1:t})$ denote the posterior distribution after observing studies $1$ through $t$. A natural way to quantify the distance between successive posteriors $p_t$ and $p_{t+1}$ is the Wasserstein-$p$ distance,
\begin{align}
\label{eq:LM}
W_p(p_t, p_{t+1}) = \left(\inf_{\gamma \in \Pi(p_t,\,p_{t+1})} \int |x_1 - x_2|^p \, d\gamma(x_1,x_2)\right)^{1/p}.
\end{align}
This metric can be interpreted as the minimal cost of transporting the mass from the previous posterior distribution $p_t$ to the updated posterior $p_{t+1}$~\citep{Villani_2008}.

The Wasserstein distance is particularly appealing for sequential meta-analysis. It is computationally convenient, especially when the posterior distributions are one-dimensional. Moreover, when the posteriors are normal, the Wasserstein-$2$ distance has a closed-form solution:
\begin{align}
W_2^2\bigl(\mathcal{N}(\mu_t,\sigma_t^2),\,\mathcal{N}(\mu_{t+1},\sigma_{t+1}^2)\bigr) = (\mu_{t+1} - \mu_t)^2 + (\sigma_{t+1} - \sigma_t)^2.\nonumber
\end{align}
In this expression, shifts in the posterior mean ($\mu$) and changes in the posterior standard deviation ($\sigma$) both contribute to the measure. The Wasserstein distance metric thus captures the possibility that an important study might not only refine our estimate but could also substantially change our uncertainty about the target parameter~\citep{mikhaeil2025b}.



\subsection{Tracking scientific consensus}

We have described an approach to sequential meta-analysis that involves three ingredients: (1) a Bayesian meta-analysis model, which formalizes assumptions about study biases and the target parameter; (2) a method for sequential Bayesian updating, derived from the model, which revises our collective beliefs about the target parameter after each new study; and
(3) a metric of posterior change, which quantifies how much each study contributes to scientific understanding. Together, these ingredients allow us to systematically analyze sequences of studies, explicitly track how scientific consensus emerges or shifts over time, and determine which studies were most important in forming that consensus. Next, we illustrate this framework using synthetic examples and real datasets.
